% !TEX root = ../LosningsforslagTRES0312.tex
\chapter{Statikk}
\label{LF:Chapter:statikk}
Dette kapittelet inneholder løsningsforslag til oppgavene i Kapittel~\ref{komp-Chapter:statikk}, \textit{Statikk}, i kompendiet.

\oppgaver{Lampe i to snorer}{k6oppg1}{%
%
En lampe med masse $m = \SI{4.0}{\kilo\gram}$ henger i ro i to snorer. Snor 1 går horisontalt inn til en vegg. Snor 2 går skrått opp til taket og danner vinkelen $\alpha = 40^\circ$ med horisontalen.
\begin{enumerate}[a)]
    \item Tegn kreftene som virker på lampa.
    \item Finn snorkraften $S_2$ i den skrå snora. (Hint: se på kreftene i $y$-retning.)
    \item Finn snorkraften $S_1$ i den horisontale snora.
\end{enumerate}
%
}

\svar[title={Løsningsforslag}]{%
\noindent
\textbf{a)} Tre krefter virker på lampa: tyngdekraften $\vec{G}$ rett ned, snorkraften $\vec{S}_1$ horisontalt inn mot veggen, og snorkraften $\vec{S}_2$ skrått oppover langs snor 2, i vinkelen $\alpha$ med horisontalen. Vi velger $x$-aksen horisontalt bort fra veggen og $y$-aksen vertikalt oppover.

\smallskip\noindent
\textbf{b)} Lampa er i ro, så $\sum F_y = 0$. Bare $S_2$ har en komponent i $y$-retning:
\[
    S_2\sin\alpha - m\mathrm{g} = 0
    \quad\implies\quad
    S_2 = \frac{m\mathrm{g}}{\sin\alpha} = \frac{\SI{4.0}{\kilo\gram}\cdot\SI{9.81}{\metre\per\second\squared}}{\sin 40^\circ} = \doubleunderline{\SI{61}{\newton}}
\]

\smallskip\noindent
\textbf{c)} $\sum F_x = 0$ gir at den horisontale komponenten av $S_2$ må oppveie $S_1$:
\[
    S_1 = S_2\cos\alpha = \SI{61.0}{\newton}\cdot\cos 40^\circ = \doubleunderline{\SI{47}{\newton}}
\]
}

\oppgaver{Dørhengsler}{k6oppg2}{%
%
%
En dør er \SI{90}{\centi\metre} bred. Du trykker vinkelrett på døren med en kraft på \SI{15}{\newton}.
\begin{enumerate}[a)]
  \item Hva er kraftmomentet om hengslene når du trykker på kanten av døren (dvs.\ \SI{90}{\centi\metre} fra hengslene)?
  \item Hva er kraftmomentet dersom du i stedet trykker på midten av døren (\SI{45}{\centi\metre} fra hengslene)?
  \item Hva skjer med kraften som trengs for å oppnå det samme momentet dersom armen halveres?
\end{enumerate}
%
}

\svar[title={Løsningsforslag}]{%
\noindent
Kraften står vinkelrett på døra, så momentarmen er avstanden fra hengslene, og $\tau = d\cdot F$ (se \textit{Kraftmoment} i Kapittel~\ref{komp-Chapter:statikk}).

\smallskip\noindent
\textbf{a)}
\[
    \tau = \SI{0.90}{\metre}\cdot\SI{15}{\newton} = \SI{13.5}{\newton\metre} \approx \doubleunderline{\SI{14}{\newton\metre}}
\]

\smallskip\noindent
\textbf{b)}
\[
    \tau = \SI{0.45}{\metre}\cdot\SI{15}{\newton} = \SI{6.75}{\newton\metre} \approx \doubleunderline{\SI{6.8}{\newton\metre}}
\]

\smallskip\noindent
\textbf{c)} Siden $\tau = d\cdot F$, må $F$ \doubleunderline{dobles} når $d$ halveres for at momentet skal bli det samme. Det er derfor det er tungt å åpne en dør ved å dytte nær hengslene.
}

\oppgaver{Stige mot vegg}{k6oppg3}{%
%
%
En ensartet stige med masse $m = \SI{15}{\kilo\gram}$ og lengde $L = \SI{4.0}{\metre}$ hviler mot en glatt vegg (ingen friksjon) i en vinkel $\theta = 65°$ med horisontal. Gulvet har friksjon. En person med masse $M = \SI{70}{\kilo\gram}$ klatrer \SI{2.5}{\metre} opp stigen.
\begin{enumerate}[a)]
  \item Finn normalkraften fra veggen.
  \item Finn normalkraften og friksjonskraften fra gulvet.
\end{enumerate}
\textit{Hint: Velg dreiepunkt i bunnen av stigen — da faller normalkraften og friksjonskraften fra gulvet ut av momentligningen.}
%
}

\svar[title={Løsningsforslag}]{%
\noindent
\textbf{Krefter:} normalkraft fra veggen $N_v$ (horisontal, fordi veggen er glatt), normalkraft fra gulvet $N_g$ (vertikal), friksjon fra gulvet $f$ (horisontal), tyngdekraften på stigen $m\mathrm{g}$ (angriper midt på stigen) og tyngdekraften på personen $M\mathrm{g}$ (angriper $d = \SI{2.5}{\metre}$ opp langs stigen).

\smallskip\noindent
\textbf{a)} Vi tar momenter om bunnen av stigen, slik at $N_g$ og $f$ faller ut. Momentarmen til en vertikal kraft er den horisontale avstanden, og momentarmen til en horisontal kraft er den vertikale avstanden. Med positiv dreieretning mot klokken:
\[
    N_v\cdot L\sin\theta - m\mathrm{g}\cdot\frac{L}{2}\cos\theta - M\mathrm{g}\cdot d\cos\theta = 0
\]
Vi løser for $N_v$:
\[
    N_v = \left(\frac{m}{2} + \frac{Md}{L}\right)\mathrm{g}\,\frac{\cos\theta}{\sin\theta}
    = \left(\SI{7.5}{\kilo\gram} + \SI{43.75}{\kilo\gram}\right)\cdot\SI{9.81}{\metre\per\second\squared}\cdot\frac{\cos 65^\circ}{\sin 65^\circ}
    = \SI{234}{\newton} \approx \doubleunderline{\SI{2.3e2}{\newton}}
\]

\smallskip\noindent
\textbf{b)} Kraftlikevekt vertikalt:
\[
    N_g = (m + M)\,\mathrm{g} = \SI{85}{\kilo\gram}\cdot\SI{9.81}{\metre\per\second\squared} = \SI{834}{\newton} \approx \doubleunderline{\SI{8.3e2}{\newton}}
\]
Kraftlikevekt horisontalt: friksjonen må oppveie normalkraften fra veggen,
\[
    f = N_v = \SI{234}{\newton} \approx \doubleunderline{\SI{2.3e2}{\newton}}
\]
}

\oppgaver{Tre lodd på en stav}{k6oppg4}{%
%
Tre lodd henger på en lett stav (staven har så liten masse at vi kan se bort fra den). Et lodd med masse $\SI{1.0}{\kilo\gram}$ henger i den venstre enden ($x = 0$), et lodd med masse $\SI{2.0}{\kilo\gram}$ henger i $x = \SI{0.50}{\metre}$, og et lodd med masse $\SI{3.0}{\kilo\gram}$ henger i $x = \SI{1.20}{\metre}$.
\begin{enumerate}[a)]
    \item Finn massesenteret $x_\text{ms}$ til systemet.
    \item Hvor må du holde staven for at den skal balansere?
\end{enumerate}
%
}

\svar[title={Løsningsforslag}]{%
\noindent
\textbf{a)} Vi bruker definisjonen av massesenter (se \textit{Massesenter} i Kapittel~\ref{komp-Chapter:statikk}):
\[
    x_\text{ms} = \frac{m_1x_1 + m_2x_2 + m_3x_3}{m_1 + m_2 + m_3}
    = \frac{\SI{1.0}{\kilo\gram}\cdot 0 + \SI{2.0}{\kilo\gram}\cdot\SI{0.50}{\metre} + \SI{3.0}{\kilo\gram}\cdot\SI{1.20}{\metre}}{\SI{6.0}{\kilo\gram}}
    = \doubleunderline{\SI{0.77}{\metre}}
\]

\smallskip\noindent
\textbf{b)} Staven balanserer når den støttes i massesenteret. Da blir det totale kraftmomentet fra tyngdekreftene null. Du må holde staven \doubleunderline{\SI{0.77}{\metre} fra den venstre enden}.
}

\oppgaver{Bilde på vegg}{k6oppg5}{%
%
Et bilde med masse $m = \SI{1.5}{\kilo\gram}$ henger i en snor over en spiker. De to delene av snora danner like store vinkler $\alpha$ med horisontalen.
\begin{enumerate}[a)]
    \item Finn snorkraften når $\alpha = 15^\circ$.
    \item Finn snorkraften når $\alpha = 45^\circ$.
    \item Hvorfor er det lurt å ikke stramme snora for mye når man henger opp et tungt bilde?
\end{enumerate}
%
}

\svar[title={Løsningsforslag}]{%
\noindent
Dette er samme situasjon som eksempelet \textit{Trafikklys} i Kapittel~\ref{komp-Chapter:statikk}. De horisontale komponentene opphever hverandre, og de vertikale komponentene må til sammen bære tyngden:
\[
    2S\sin\alpha = m\mathrm{g}
    \quad\implies\quad
    S = \frac{m\mathrm{g}}{2\sin\alpha}
\]

\smallskip\noindent
\textbf{a)}
\[
    S = \frac{\SI{1.5}{\kilo\gram}\cdot\SI{9.81}{\metre\per\second\squared}}{2\sin 15^\circ} = \doubleunderline{\SI{28}{\newton}}
\]

\smallskip\noindent
\textbf{b)}
\[
    S = \frac{\SI{1.5}{\kilo\gram}\cdot\SI{9.81}{\metre\per\second\squared}}{2\sin 45^\circ} = \doubleunderline{\SI{10}{\newton}}
\]

\smallskip\noindent
\textbf{c)} En stram snor gir en liten vinkel $\alpha$, og da blir $\sin\alpha$ liten og snorkraften stor. Med $\alpha = 15^\circ$ er snorkraften nesten dobbelt så stor som tyngden til bildet. En for stram snor kan derfor ryke, eller rive spikeren ut av veggen.
}

\oppgaver{Vippehuske}{k6oppg6}{%
%
En vippehuske er $\SI{4.0}{\metre}$ lang og har støttepunktet på midten. Se bort fra massen til selve vippa. Et barn med masse $\SI{25}{\kilo\gram}$ sitter helt ytterst på den ene siden.
\begin{enumerate}[a)]
    \item Hvor langt fra støttepunktet må en voksen med masse $\SI{80}{\kilo\gram}$ sitte for at vippa skal balansere?
    \item Hvor stor kraft virker det fra støttepunktet på vippa når den balanserer?
\end{enumerate}
%
}

\svar[title={Løsningsforslag}]{%
\noindent
\textbf{a)} Vi tar momenter om støttepunktet. Barnet sitter $\SI{2.0}{\metre}$ fra støttepunktet. For at vippa skal balansere, må momentene på hver side være like store:
\[
    m_\text{barn}\mathrm{g}\cdot d_\text{barn} = m_\text{voksen}\mathrm{g}\cdot d_\text{voksen}
    \quad\implies\quad
    d_\text{voksen} = \frac{\SI{25}{\kilo\gram}\cdot\SI{2.0}{\metre}}{\SI{80}{\kilo\gram}} = \doubleunderline{\SI{0.63}{\metre}}
\]

\smallskip\noindent
\textbf{b)} Kraftlikevekt vertikalt: kraften fra støttepunktet må bære begge personene:
\[
    N = (\SI{25}{\kilo\gram} + \SI{80}{\kilo\gram})\cdot\SI{9.81}{\metre\per\second\squared} = \SI{1.03}{\kilo\newton} \approx \doubleunderline{\SI{1.0}{\kilo\newton}}
\]
}

\oppgaver{To personer bærer en planke}{k6oppg7}{%
%
To personer bærer en ensartet planke som er $\SI{3.0}{\metre}$ lang og har masse $\SI{30}{\kilo\gram}$. Person A holder i den venstre enden, og person B holder $\SI{0.50}{\metre}$ fra den høyre enden. En kasse med masse $\SI{20}{\kilo\gram}$ står på planken, $\SI{1.0}{\metre}$ fra den venstre enden.

Finn hvor stor kraft hver av personene må løfte med. (Hint: velg dreiepunkt der person A holder.)
%
}

\svar[title={Løsningsforslag}]{%
\noindent
Vi måler avstander fra den venstre enden, der person A holder. Tyngdekraften på planken angriper midt på, $\SI{1.5}{\metre}$ fra A. Kassen står $\SI{1.0}{\metre}$ fra A, og person B holder $\SI{3.0}{\metre} - \SI{0.50}{\metre} = \SI{2.5}{\metre}$ fra A.

\smallskip\noindent
\textbf{Momentlikning om A} (positiv dreieretning mot klokken):
\[
    F_B\cdot\SI{2.5}{\metre} - \SI{30}{\kilo\gram}\cdot\mathrm{g}\cdot\SI{1.5}{\metre} - \SI{20}{\kilo\gram}\cdot\mathrm{g}\cdot\SI{1.0}{\metre} = 0
\]
\[
    F_B = \frac{(\SI{45}{\kilo\gram\metre} + \SI{20}{\kilo\gram\metre})\cdot\SI{9.81}{\metre\per\second\squared}}{\SI{2.5}{\metre}} = \SI{255}{\newton} \approx \doubleunderline{\SI{2.6e2}{\newton}}
\]

\smallskip\noindent
\textbf{Kraftlikevekt vertikalt:}
\[
    F_A = (\SI{30}{\kilo\gram} + \SI{20}{\kilo\gram})\cdot\mathrm{g} - F_B = \SI{490.5}{\newton} - \SI{255}{\newton} = \SI{235}{\newton} \approx \doubleunderline{\SI{2.4e2}{\newton}}
\]
Person B bærer litt mer enn A. B holder nærmere midten av planken enn A, men kassen står nærmest A, så forskjellen blir liten.
}

\oppgaver{Massesenteret til jorda og månen}{k6oppg8}{%
%
Jorda har masse $\SI{5.97e24}{\kilo\gram}$ og radius $\SI{6.37e6}{\metre}$. Månen har masse $\SI{7.35e22}{\kilo\gram}$, og avstanden mellom sentrene til jorda og månen er $\SI{3.84e8}{\metre}$.
\begin{enumerate}[a)]
    \item Finn massesenteret til jord--måne-systemet, målt fra jordas sentrum.
    \item Ligger massesenteret inne i jorda eller utenfor?
\end{enumerate}
%
}

\svar[title={Løsningsforslag}]{%
\noindent
\textbf{a)} Vi legger $x = 0$ i sentrum av jorda og $x = d = \SI{3.84e8}{\metre}$ i sentrum av månen:
\[
    x_\text{ms} = \frac{m_J\cdot 0 + m_M\cdot d}{m_J + m_M}
    = \frac{\SI{7.35e22}{\kilo\gram}\cdot\SI{3.84e8}{\metre}}{\SI{5.97e24}{\kilo\gram} + \SI{7.35e22}{\kilo\gram}}
    = \doubleunderline{\SI{4.67e6}{\metre}}
\]

\smallskip\noindent
\textbf{b)} Jordas radius er $\SI{6.37e6}{\metre}$, så massesenteret ligger \doubleunderline{inne i jorda}, omtrent \SI{1700}{\kilo\metre} under overflaten. Det er rundt dette punktet både jorda og månen går i bane.
}

\oppgaver{Mutter som sitter fast}{k6oppg9}{%
%
En mutter sitter fast, og det trengs et kraftmoment på $\SI{50}{\newton\metre}$ for å løsne den. Du har en skiftenøkkel som er $\SI{30}{\centi\metre}$ lang.
\begin{enumerate}[a)]
    \item Hvor stor kraft må du minst bruke, og i hvilken retning bør du dra?
    \item Hvor stor kraft trengs dersom kraften danner en vinkel på $60^\circ$ med skaftet på nøkkelen?
    \item Du setter et rør over skaftet slik at armen blir $\SI{60}{\centi\metre}$ lang. Hvor stor kraft trengs nå, når du drar vinkelrett?
\end{enumerate}
%
}

\svar[title={Løsningsforslag}]{%
\noindent
Vi bruker $\tau = r\cdot F\cdot\sin\theta$, der $\theta$ er vinkelen mellom skaftet og kraften.

\smallskip\noindent
\textbf{a)} Kraften blir minst når $\sin\theta$ er størst, altså når du drar \doubleunderline{vinkelrett på skaftet} ($\theta = 90^\circ$):
\[
    F = \frac{\tau}{r} = \frac{\SI{50}{\newton\metre}}{\SI{0.30}{\metre}} = \SI{167}{\newton} \approx \doubleunderline{\SI{1.7e2}{\newton}}
\]

\smallskip\noindent
\textbf{b)}
\[
    F = \frac{\tau}{r\sin 60^\circ} = \frac{\SI{50}{\newton\metre}}{\SI{0.30}{\metre}\cdot 0.866} = \SI{192}{\newton} \approx \doubleunderline{\SI{1.9e2}{\newton}}
\]

\smallskip\noindent
\textbf{c)}
\[
    F = \frac{\tau}{r} = \frac{\SI{50}{\newton\metre}}{\SI{0.60}{\metre}} = \doubleunderline{\SI{83}{\newton}}
\]
Dobbelt så lang arm gir halvparten så stor kraft.
}
